% !TEX root = ../../main.tex
%% Défi du chapitre 20 — Le point de Fermat
%% Résultats vérifiés par verifications/defis_16_20.py ;
%% lignes de niveau générées par verifications/fig_fermat_niveaux.py.

\rubrique{Défi}

\begin{defi}{Le point de Fermat}{2 h}
Soient $A$, $B$, $C$ trois points non alignés du plan $\R^{2}$, muni de la
norme euclidienne. On cherche le point $M$ qui minimise la somme des
distances $f(M)=MA+MB+MC$ — question posée par Pierre de Fermat à
Evangelista Torricelli vers 1640, et problème concret : où relier trois
villes par un réseau de longueur minimale ?

\partieExo{A}{Gradient et points critiques}
\begin{questions}
  \item Pour $P\in\R^{2}$ fixé, on pose $d_P(M)=\norme{M-P}$. Montrer que
    $d_P$ est de classe $\mathcal{C}^{1}$ sur $\R^{2}\setminus\ens{P}$, de
    gradient le vecteur \emph{unitaire}
    $\nabla d_P(M)=\dfrac{M-P}{\norme{M-P}}$.
  \item En déduire que $f$ est de classe $\mathcal{C}^{1}$ sur
    $U=\R^{2}\setminus\ens{A,B,C}$, et que ses points critiques dans $U$
    sont les $M$ tels que $u_A+u_B+u_C=0$, où $u_P$ désigne le vecteur
    unitaire dirigé de $P$ vers $M$.
  \item Montrer que trois vecteurs unitaires de somme nulle font deux à deux
    des angles de $\frac{2\pi}{3}$. Conclure : en un point critique $M$, les
    angles $\widehat{AMB}$, $\widehat{BMC}$, $\widehat{CMA}$ valent tous
    $\frac{2\pi}{3}$.
\end{questions}

\partieExo{B}{Existence et position du minimum}
\begin{questions}[resume]
  \item Montrer que $f(M)\geqslant3\norme{M}-\bigl(\norme A+\norme B+\norme C\bigr)$.
    On admet qu'une fonction continue sur un disque fermé y atteint son
    minimum ; en déduire que $f$ admet un minimum global sur $\R^{2}$.
  \item On suppose que tous les angles du triangle sont strictement
    inférieurs à $\frac{2\pi}{3}$. Pour $u$ unitaire, calculer la dérivée à
    droite en $0$ de $t\mapsto f(A+tu)$, puis choisir $u$ pour qu'elle soit
    strictement négative : le minimum n'est pas atteint en $A$, ni (de
    même) en $B$ ou $C$. En déduire qu'il est atteint en un point critique
    $F$ de $U$ : le \emph{point de Fermat}, d'où l'on voit les trois côtés
    sous des angles de $\frac{2\pi}{3}$.
\end{questions}

\medskip
\begin{center}
\begin{tikzpicture}[scale=1.25,line cap=round,line join=round]
  \begin{scope}
    \clip (-2.2,-1.9) rectangle (2.2,3.0);
    % Fichier généré par verifications/fig_fermat_niveaux.py — ne pas éditer
\draw[otNavy!45,line width=0.5pt] (-0.093,0.188) -- (-0.137,0.193) -- (-0.181,0.200) -- (-0.225,0.210) -- (-0.269,0.223) -- (-0.304,0.236) -- (-0.339,0.253) -- (-0.375,0.276) -- (-0.401,0.298) -- (-0.428,0.329) -- (-0.445,0.361) -- (-0.460,0.404) -- (-0.465,0.449) -- (-0.463,0.504) -- (-0.456,0.549) -- (-0.445,0.593) -- (-0.433,0.630) -- (-0.419,0.668) -- (-0.401,0.708) -- (-0.384,0.743) -- (-0.366,0.774) -- (-0.343,0.810) -- (-0.322,0.841) -- (-0.297,0.873) -- (-0.274,0.900) -- (-0.249,0.927) -- (-0.220,0.955) -- (-0.190,0.979) -- (-0.160,1.000) -- (-0.128,1.018) -- (-0.085,1.036) -- (-0.048,1.045) -- (0.004,1.050) -- (0.052,1.045) -- (0.093,1.033) -- (0.128,1.018) -- (0.163,0.997) -- (0.198,0.973) -- (0.225,0.950) -- (0.251,0.925) -- (0.278,0.896) -- (0.304,0.865) -- (0.325,0.837) -- (0.348,0.803) -- (0.366,0.774) -- (0.386,0.738) -- (0.404,0.702) -- (0.420,0.666) -- (0.436,0.621) -- (0.448,0.585) -- (0.458,0.540) -- (0.464,0.495) -- (0.465,0.440) -- (0.457,0.395) -- (0.445,0.359) -- (0.423,0.323) -- (0.399,0.296) -- (0.366,0.269) -- (0.335,0.251) -- (0.296,0.233) -- (0.260,0.220) -- (0.216,0.208) -- (0.172,0.198) -- (0.128,0.191) -- (0.084,0.187) -- (0.031,0.184) -- (-0.022,0.184) -- (-0.075,0.186) -- cycle;
\draw[otNavy!45,line width=0.5pt] (-0.190,-0.164) -- (-0.234,-0.160) -- (-0.287,-0.155) -- (-0.331,-0.150) -- (-0.375,-0.144) -- (-0.420,-0.137) -- (-0.472,-0.128) -- (-0.516,-0.119) -- (-0.560,-0.110) -- (-0.599,-0.101) -- (-0.639,-0.091) -- (-0.683,-0.079) -- (-0.727,-0.066) -- (-0.763,-0.056) -- (-0.807,-0.041) -- (-0.844,-0.028) -- (-0.886,-0.013) -- (-0.921,0.002) -- (-0.955,0.026) -- (-0.963,0.071) -- (-0.958,0.125) -- (-0.953,0.170) -- (-0.946,0.215) -- (-0.938,0.260) -- (-0.930,0.305) -- (-0.920,0.350) -- (-0.909,0.395) -- (-0.898,0.440) -- (-0.886,0.483) -- (-0.875,0.522) -- (-0.860,0.567) -- (-0.848,0.603) -- (-0.833,0.644) -- (-0.818,0.684) -- (-0.804,0.720) -- (-0.789,0.756) -- (-0.772,0.795) -- (-0.754,0.834) -- (-0.736,0.870) -- (-0.719,0.905) -- (-0.701,0.939) -- (-0.683,0.973) -- (-0.662,1.009) -- (-0.641,1.045) -- (-0.622,1.075) -- (-0.601,1.108) -- (-0.578,1.143) -- (-0.558,1.171) -- (-0.533,1.205) -- (-0.512,1.234) -- (-0.489,1.263) -- (-0.463,1.296) -- (-0.439,1.324) -- (-0.416,1.351) -- (-0.392,1.378) -- (-0.366,1.406) -- (-0.339,1.433) -- (-0.312,1.460) -- (-0.282,1.487) -- (-0.251,1.514) -- (-0.225,1.535) -- (-0.193,1.559) -- (-0.163,1.579) -- (-0.128,1.600) -- (-0.093,1.617) -- (-0.052,1.631) -- (-0.004,1.638) -- (0.048,1.632) -- (0.084,1.621) -- (0.121,1.604) -- (0.154,1.585) -- (0.190,1.561) -- (0.218,1.541) -- (0.251,1.514) -- (0.278,1.491) -- (0.304,1.467) -- (0.331,1.441) -- (0.357,1.414) -- (0.384,1.387) -- (0.410,1.358) -- (0.436,1.328) -- (0.462,1.297) -- (0.484,1.270) -- (0.507,1.241) -- (0.532,1.207) -- (0.552,1.180) -- (0.577,1.144) -- (0.595,1.116) -- (0.618,1.081) -- (0.639,1.047) -- (0.657,1.017) -- (0.678,0.982) -- (0.697,0.945) -- (0.716,0.909) -- (0.735,0.873) -- (0.752,0.837) -- (0.769,0.801) -- (0.785,0.765) -- (0.800,0.729) -- (0.816,0.691) -- (0.832,0.648) -- (0.845,0.612) -- (0.860,0.568) -- (0.872,0.531) -- (0.885,0.486) -- (0.895,0.449) -- (0.907,0.404) -- (0.918,0.359) -- (0.928,0.314) -- (0.937,0.269) -- (0.945,0.224) -- (0.952,0.179) -- (0.958,0.134) -- (0.962,0.080) -- (0.957,0.029) -- (0.930,0.007) -- (0.892,-0.010) -- (0.851,-0.026) -- (0.816,-0.038) -- (0.772,-0.053) -- (0.734,-0.065) -- (0.692,-0.077) -- (0.648,-0.088) -- (0.604,-0.099) -- (0.560,-0.110) -- (0.518,-0.119) -- (0.472,-0.128) -- (0.428,-0.135) -- (0.384,-0.142) -- (0.339,-0.149) -- (0.290,-0.155) -- (0.242,-0.160) -- (0.194,-0.164) -- (0.145,-0.167) -- (0.093,-0.169) -- (0.040,-0.171) -- (-0.013,-0.171) -- (-0.066,-0.170) -- (-0.119,-0.168) -- (-0.172,-0.165) -- cycle;
\draw[otNavy!45,line width=0.5pt] (-0.207,-0.525) -- (-0.251,-0.520) -- (-0.296,-0.515) -- (-0.348,-0.508) -- (-0.392,-0.502) -- (-0.436,-0.494) -- (-0.481,-0.485) -- (-0.525,-0.476) -- (-0.569,-0.465) -- (-0.613,-0.453) -- (-0.648,-0.443) -- (-0.692,-0.429) -- (-0.731,-0.416) -- (-0.772,-0.401) -- (-0.807,-0.387) -- (-0.844,-0.371) -- (-0.882,-0.353) -- (-0.917,-0.335) -- (-0.949,-0.317) -- (-0.983,-0.296) -- (-1.018,-0.272) -- (-1.045,-0.253) -- (-1.076,-0.227) -- (-1.104,-0.200) -- (-1.129,-0.173) -- (-1.151,-0.146) -- (-1.175,-0.110) -- (-1.195,-0.076) -- (-1.212,-0.038) -- (-1.226,-0.001) -- (-1.239,0.040) -- (-1.248,0.080) -- (-1.257,0.132) -- (-1.261,0.179) -- (-1.264,0.233) -- (-1.264,0.287) -- (-1.261,0.341) -- (-1.257,0.394) -- (-1.251,0.440) -- (-1.244,0.486) -- (-1.237,0.531) -- (-1.227,0.576) -- (-1.217,0.621) -- (-1.206,0.666) -- (-1.195,0.708) -- (-1.183,0.747) -- (-1.169,0.792) -- (-1.157,0.828) -- (-1.142,0.870) -- (-1.127,0.909) -- (-1.113,0.945) -- (-1.098,0.982) -- (-1.080,1.022) -- (-1.063,1.061) -- (-1.045,1.098) -- (-1.027,1.133) -- (-1.010,1.167) -- (-0.992,1.200) -- (-0.973,1.234) -- (-0.952,1.270) -- (-0.931,1.306) -- (-0.913,1.335) -- (-0.891,1.369) -- (-0.869,1.403) -- (-0.849,1.432) -- (-0.824,1.467) -- (-0.803,1.496) -- (-0.780,1.526) -- (-0.755,1.559) -- (-0.734,1.586) -- (-0.710,1.615) -- (-0.683,1.646) -- (-0.657,1.676) -- (-0.632,1.703) -- (-0.607,1.730) -- (-0.581,1.757) -- (-0.554,1.784) -- (-0.526,1.811) -- (-0.498,1.838) -- (-0.472,1.862) -- (-0.445,1.885) -- (-0.415,1.910) -- (-0.384,1.936) -- (-0.357,1.956) -- (-0.322,1.982) -- (-0.295,2.001) -- (-0.260,2.024) -- (-0.225,2.045) -- (-0.190,2.064) -- (-0.154,2.081) -- (-0.119,2.095) -- (-0.075,2.108) -- (-0.031,2.116) -- (0.022,2.116) -- (0.072,2.109) -- (0.110,2.098) -- (0.152,2.082) -- (0.190,2.064) -- (0.224,2.046) -- (0.254,2.028) -- (0.287,2.007) -- (0.321,1.983) -- (0.348,1.963) -- (0.382,1.937) -- (0.410,1.915) -- (0.437,1.892) -- (0.468,1.865) -- (0.498,1.838) -- (0.525,1.813) -- (0.551,1.787) -- (0.578,1.761) -- (0.604,1.733) -- (0.630,1.705) -- (0.657,1.676) -- (0.681,1.649) -- (0.704,1.622) -- (0.727,1.593) -- (0.754,1.560) -- (0.776,1.532) -- (0.798,1.503) -- (0.823,1.469) -- (0.842,1.441) -- (0.867,1.405) -- (0.886,1.377) -- (0.908,1.342) -- (0.930,1.307) -- (0.948,1.278) -- (0.968,1.243) -- (0.988,1.207) -- (1.008,1.171) -- (1.026,1.135) -- (1.044,1.099) -- (1.062,1.063) -- (1.078,1.027) -- (1.094,0.991) -- (1.109,0.955) -- (1.124,0.917) -- (1.141,0.873) -- (1.154,0.837) -- (1.168,0.794) -- (1.180,0.756) -- (1.194,0.711) -- (1.204,0.675) -- (1.215,0.630) -- (1.226,0.585) -- (1.235,0.540) -- (1.243,0.495) -- (1.250,0.449) -- (1.256,0.395) -- (1.260,0.350) -- (1.263,0.296) -- (1.264,0.242) -- (1.262,0.188) -- (1.257,0.134) -- (1.250,0.089) -- (1.240,0.044) -- (1.229,0.008) -- (1.213,-0.037) -- (1.196,-0.074) -- (1.177,-0.107) -- (1.157,-0.137) -- (1.133,-0.168) -- (1.107,-0.197) -- (1.080,-0.223) -- (1.054,-0.245) -- (1.019,-0.272) -- (0.992,-0.291) -- (0.957,-0.313) -- (0.921,-0.333) -- (0.886,-0.351) -- (0.851,-0.368) -- (0.816,-0.383) -- (0.779,-0.398) -- (0.736,-0.414) -- (0.701,-0.426) -- (0.657,-0.440) -- (0.616,-0.452) -- (0.578,-0.463) -- (0.533,-0.473) -- (0.489,-0.483) -- (0.445,-0.492) -- (0.401,-0.500) -- (0.357,-0.507) -- (0.304,-0.514) -- (0.260,-0.520) -- (0.209,-0.524) -- (0.163,-0.528) -- (0.110,-0.531) -- (0.057,-0.533) -- (0.004,-0.533) -- (-0.048,-0.533) -- (-0.101,-0.531) -- (-0.154,-0.528) -- (-0.207,-0.525) -- cycle;
\draw[otNavy!45,line width=0.5pt] (-0.057,-0.903) -- (-0.101,-0.902) -- (-0.154,-0.899) -- (-0.207,-0.895) -- (-0.251,-0.890) -- (-0.295,-0.885) -- (-0.348,-0.877) -- (-0.392,-0.870) -- (-0.436,-0.862) -- (-0.481,-0.853) -- (-0.525,-0.843) -- (-0.569,-0.831) -- (-0.604,-0.822) -- (-0.648,-0.809) -- (-0.690,-0.795) -- (-0.727,-0.782) -- (-0.765,-0.768) -- (-0.807,-0.751) -- (-0.842,-0.736) -- (-0.877,-0.720) -- (-0.913,-0.703) -- (-0.948,-0.685) -- (-0.983,-0.666) -- (-1.018,-0.646) -- (-1.054,-0.624) -- (-1.081,-0.606) -- (-1.115,-0.582) -- (-1.144,-0.561) -- (-1.177,-0.535) -- (-1.204,-0.513) -- (-1.231,-0.488) -- (-1.259,-0.461) -- (-1.286,-0.434) -- (-1.311,-0.407) -- (-1.336,-0.378) -- (-1.362,-0.345) -- (-1.383,-0.317) -- (-1.406,-0.283) -- (-1.425,-0.254) -- (-1.446,-0.218) -- (-1.465,-0.182) -- (-1.483,-0.146) -- (-1.499,-0.110) -- (-1.513,-0.074) -- (-1.529,-0.028) -- (-1.541,0.008) -- (-1.553,0.053) -- (-1.564,0.098) -- (-1.573,0.143) -- (-1.580,0.188) -- (-1.585,0.233) -- (-1.590,0.287) -- (-1.593,0.332) -- (-1.594,0.386) -- (-1.594,0.440) -- (-1.592,0.493) -- (-1.588,0.540) -- (-1.583,0.594) -- (-1.577,0.639) -- (-1.571,0.684) -- (-1.563,0.729) -- (-1.554,0.774) -- (-1.544,0.819) -- (-1.533,0.864) -- (-1.521,0.909) -- (-1.511,0.945) -- (-1.497,0.991) -- (-1.485,1.027) -- (-1.469,1.072) -- (-1.456,1.108) -- (-1.442,1.145) -- (-1.424,1.188) -- (-1.408,1.225) -- (-1.392,1.261) -- (-1.375,1.297) -- (-1.357,1.333) -- (-1.338,1.369) -- (-1.319,1.405) -- (-1.301,1.439) -- (-1.283,1.470) -- (-1.262,1.505) -- (-1.240,1.541) -- (-1.221,1.570) -- (-1.199,1.604) -- (-1.177,1.635) -- (-1.155,1.667) -- (-1.133,1.696) -- (-1.108,1.730) -- (-1.086,1.757) -- (-1.063,1.786) -- (-1.036,1.818) -- (-1.011,1.847) -- (-0.986,1.874) -- (-0.961,1.901) -- (-0.935,1.928) -- (-0.909,1.956) -- (-0.881,1.983) -- (-0.852,2.010) -- (-0.824,2.035) -- (-0.798,2.058) -- (-0.769,2.082) -- (-0.736,2.108) -- (-0.710,2.129) -- (-0.676,2.154) -- (-0.648,2.173) -- (-0.613,2.197) -- (-0.581,2.217) -- (-0.551,2.235) -- (-0.516,2.255) -- (-0.481,2.274) -- (-0.445,2.292) -- (-0.410,2.308) -- (-0.369,2.325) -- (-0.331,2.340) -- (-0.295,2.352) -- (-0.251,2.365) -- (-0.207,2.376) -- (-0.163,2.385) -- (-0.119,2.392) -- (-0.068,2.397) -- (-0.022,2.400) -- (0.031,2.399) -- (0.075,2.397) -- (0.128,2.391) -- (0.172,2.384) -- (0.216,2.374) -- (0.260,2.363) -- (0.295,2.352) -- (0.339,2.337) -- (0.375,2.323) -- (0.412,2.307) -- (0.451,2.289) -- (0.487,2.271) -- (0.520,2.253) -- (0.552,2.235) -- (0.586,2.214) -- (0.622,2.191) -- (0.650,2.172) -- (0.683,2.148) -- (0.712,2.127) -- (0.745,2.101) -- (0.772,2.080) -- (0.801,2.055) -- (0.832,2.028) -- (0.860,2.002) -- (0.886,1.977) -- (0.913,1.952) -- (0.939,1.925) -- (0.966,1.897) -- (0.992,1.868) -- (1.018,1.839) -- (1.042,1.811) -- (1.064,1.784) -- (1.089,1.754) -- (1.114,1.721) -- (1.135,1.694) -- (1.160,1.660) -- (1.180,1.631) -- (1.204,1.596) -- (1.223,1.568) -- (1.246,1.532) -- (1.265,1.499) -- (1.284,1.469) -- (1.304,1.432) -- (1.324,1.396) -- (1.343,1.360) -- (1.361,1.324) -- (1.379,1.288) -- (1.396,1.252) -- (1.412,1.216) -- (1.427,1.180) -- (1.442,1.144) -- (1.459,1.099) -- (1.473,1.063) -- (1.486,1.025) -- (1.500,0.982) -- (1.512,0.941) -- (1.524,0.900) -- (1.535,0.855) -- (1.546,0.810) -- (1.556,0.765) -- (1.565,0.720) -- (1.572,0.675) -- (1.579,0.630) -- (1.584,0.585) -- (1.589,0.531) -- (1.592,0.486) -- (1.594,0.431) -- (1.594,0.377) -- (1.593,0.323) -- (1.590,0.278) -- (1.584,0.224) -- (1.578,0.179) -- (1.571,0.134) -- (1.562,0.089) -- (1.551,0.044) -- (1.539,0.001) -- (1.526,-0.037) -- (1.512,-0.076) -- (1.495,-0.119) -- (1.479,-0.155) -- (1.461,-0.191) -- (1.442,-0.226) -- (1.424,-0.256) -- (1.402,-0.290) -- (1.380,-0.322) -- (1.356,-0.353) -- (1.334,-0.380) -- (1.309,-0.409) -- (1.283,-0.437) -- (1.257,-0.464) -- (1.230,-0.489) -- (1.200,-0.515) -- (1.168,-0.542) -- (1.142,-0.562) -- (1.108,-0.588) -- (1.080,-0.606) -- (1.045,-0.629) -- (1.010,-0.651) -- (0.978,-0.669) -- (0.945,-0.687) -- (0.910,-0.705) -- (0.872,-0.723) -- (0.833,-0.740) -- (0.798,-0.755) -- (0.763,-0.769) -- (0.719,-0.785) -- (0.683,-0.797) -- (0.639,-0.811) -- (0.603,-0.822) -- (0.560,-0.834) -- (0.516,-0.845) -- (0.472,-0.855) -- (0.428,-0.864) -- (0.384,-0.872) -- (0.339,-0.879) -- (0.294,-0.885) -- (0.242,-0.891) -- (0.198,-0.895) -- (0.145,-0.899) -- (0.093,-0.902) -- (0.048,-0.903) -- (-0.004,-0.904) -- (-0.057,-0.903) -- cycle;
\draw[otNavy!45,line width=0.5pt] (-0.192,-1.300) -- (-0.242,-1.295) -- (-0.287,-1.291) -- (-0.339,-1.284) -- (-0.384,-1.277) -- (-0.428,-1.269) -- (-0.472,-1.261) -- (-0.516,-1.252) -- (-0.560,-1.241) -- (-0.604,-1.230) -- (-0.645,-1.219) -- (-0.683,-1.208) -- (-0.727,-1.194) -- (-0.763,-1.182) -- (-0.807,-1.166) -- (-0.842,-1.153) -- (-0.880,-1.138) -- (-0.921,-1.120) -- (-0.957,-1.104) -- (-0.992,-1.088) -- (-1.027,-1.070) -- (-1.063,-1.052) -- (-1.098,-1.032) -- (-1.133,-1.012) -- (-1.164,-0.993) -- (-1.195,-0.974) -- (-1.230,-0.951) -- (-1.260,-0.930) -- (-1.292,-0.908) -- (-1.322,-0.885) -- (-1.354,-0.860) -- (-1.380,-0.838) -- (-1.410,-0.813) -- (-1.440,-0.786) -- (-1.468,-0.759) -- (-1.495,-0.733) -- (-1.521,-0.706) -- (-1.547,-0.678) -- (-1.571,-0.651) -- (-1.594,-0.624) -- (-1.618,-0.595) -- (-1.644,-0.561) -- (-1.664,-0.533) -- (-1.689,-0.498) -- (-1.707,-0.470) -- (-1.730,-0.434) -- (-1.750,-0.400) -- (-1.768,-0.369) -- (-1.786,-0.335) -- (-1.804,-0.299) -- (-1.821,-0.263) -- (-1.838,-0.223) -- (-1.856,-0.182) -- (-1.869,-0.146) -- (-1.883,-0.108) -- (-1.897,-0.065) -- (-1.909,-0.023) -- (-1.920,0.017) -- (-1.931,0.062) -- (-1.940,0.107) -- (-1.949,0.152) -- (-1.956,0.197) -- (-1.963,0.242) -- (-1.969,0.296) -- (-1.973,0.341) -- (-1.976,0.395) -- (-1.977,0.449) -- (-1.978,0.504) -- (-1.977,0.558) -- (-1.974,0.612) -- (-1.971,0.659) -- (-1.966,0.711) -- (-1.961,0.756) -- (-1.953,0.810) -- (-1.946,0.855) -- (-1.937,0.900) -- (-1.928,0.945) -- (-1.918,0.991) -- (-1.909,1.028) -- (-1.898,1.072) -- (-1.885,1.117) -- (-1.874,1.154) -- (-1.860,1.198) -- (-1.847,1.234) -- (-1.831,1.279) -- (-1.817,1.315) -- (-1.803,1.351) -- (-1.786,1.393) -- (-1.768,1.432) -- (-1.752,1.469) -- (-1.735,1.505) -- (-1.717,1.541) -- (-1.698,1.577) -- (-1.680,1.611) -- (-1.662,1.643) -- (-1.643,1.676) -- (-1.622,1.712) -- (-1.600,1.747) -- (-1.582,1.775) -- (-1.559,1.811) -- (-1.539,1.840) -- (-1.515,1.874) -- (-1.495,1.902) -- (-1.468,1.937) -- (-1.447,1.965) -- (-1.424,1.993) -- (-1.398,2.025) -- (-1.372,2.055) -- (-1.348,2.082) -- (-1.324,2.109) -- (-1.299,2.136) -- (-1.272,2.163) -- (-1.245,2.190) -- (-1.217,2.217) -- (-1.189,2.244) -- (-1.160,2.270) -- (-1.133,2.294) -- (-1.106,2.316) -- (-1.073,2.343) -- (-1.045,2.366) -- (-1.015,2.388) -- (-0.983,2.412) -- (-0.952,2.433) -- (-0.921,2.454) -- (-0.886,2.477) -- (-0.855,2.497) -- (-0.824,2.515) -- (-0.789,2.535) -- (-0.754,2.554) -- (-0.719,2.572) -- (-0.683,2.590) -- (-0.648,2.606) -- (-0.609,2.623) -- (-0.569,2.639) -- (-0.533,2.653) -- (-0.489,2.668) -- (-0.454,2.679) -- (-0.410,2.692) -- (-0.366,2.704) -- (-0.327,2.713) -- (-0.284,2.722) -- (-0.234,2.731) -- (-0.190,2.737) -- (-0.145,2.742) -- (-0.093,2.746) -- (-0.040,2.749) -- (0.004,2.749) -- (0.048,2.748) -- (0.101,2.746) -- (0.154,2.741) -- (0.198,2.736) -- (0.242,2.729) -- (0.287,2.721) -- (0.331,2.712) -- (0.375,2.702) -- (0.419,2.690) -- (0.462,2.677) -- (0.498,2.665) -- (0.541,2.650) -- (0.578,2.636) -- (0.613,2.621) -- (0.650,2.605) -- (0.689,2.587) -- (0.726,2.569) -- (0.760,2.551) -- (0.793,2.533) -- (0.825,2.515) -- (0.860,2.494) -- (0.895,2.472) -- (0.926,2.452) -- (0.957,2.430) -- (0.990,2.406) -- (1.018,2.386) -- (1.050,2.361) -- (1.080,2.338) -- (1.107,2.316) -- (1.138,2.289) -- (1.168,2.262) -- (1.195,2.238) -- (1.221,2.213) -- (1.248,2.188) -- (1.274,2.161) -- (1.301,2.134) -- (1.327,2.105) -- (1.354,2.076) -- (1.380,2.046) -- (1.403,2.019) -- (1.425,1.992) -- (1.451,1.960) -- (1.475,1.928) -- (1.495,1.901) -- (1.521,1.865) -- (1.540,1.838) -- (1.565,1.802) -- (1.583,1.774) -- (1.605,1.739) -- (1.627,1.704) -- (1.644,1.674) -- (1.664,1.640) -- (1.684,1.604) -- (1.703,1.568) -- (1.721,1.532) -- (1.739,1.496) -- (1.756,1.460) -- (1.772,1.423) -- (1.788,1.387) -- (1.803,1.351) -- (1.821,1.306) -- (1.834,1.270) -- (1.847,1.234) -- (1.863,1.189) -- (1.874,1.153) -- (1.888,1.108) -- (1.900,1.063) -- (1.909,1.027) -- (1.920,0.982) -- (1.930,0.936) -- (1.939,0.891) -- (1.947,0.846) -- (1.954,0.801) -- (1.962,0.747) -- (1.967,0.702) -- (1.971,0.657) -- (1.975,0.603) -- (1.977,0.549) -- (1.978,0.495) -- (1.977,0.440) -- (1.975,0.386) -- (1.972,0.332) -- (1.968,0.287) -- (1.962,0.236) -- (1.955,0.188) -- (1.947,0.143) -- (1.939,0.098) -- (1.929,0.053) -- (1.918,0.009) -- (1.907,-0.028) -- (1.894,-0.074) -- (1.882,-0.110) -- (1.866,-0.155) -- (1.852,-0.191) -- (1.837,-0.227) -- (1.821,-0.263) -- (1.803,-0.301) -- (1.786,-0.336) -- (1.767,-0.371) -- (1.746,-0.407) -- (1.724,-0.443) -- (1.706,-0.472) -- (1.683,-0.506) -- (1.662,-0.536) -- (1.637,-0.570) -- (1.616,-0.597) -- (1.592,-0.627) -- (1.565,-0.658) -- (1.539,-0.687) -- (1.514,-0.714) -- (1.487,-0.741) -- (1.459,-0.768) -- (1.433,-0.792) -- (1.406,-0.816) -- (1.378,-0.840) -- (1.345,-0.867) -- (1.318,-0.888) -- (1.285,-0.912) -- (1.257,-0.933) -- (1.221,-0.957) -- (1.193,-0.975) -- (1.160,-0.996) -- (1.124,-1.017) -- (1.089,-1.037) -- (1.054,-1.056) -- (1.018,-1.074) -- (0.983,-1.092) -- (0.948,-1.108) -- (0.913,-1.124) -- (0.877,-1.139) -- (0.835,-1.156) -- (0.798,-1.169) -- (0.760,-1.183) -- (0.719,-1.197) -- (0.676,-1.210) -- (0.639,-1.221) -- (0.595,-1.233) -- (0.551,-1.243) -- (0.507,-1.254) -- (0.463,-1.263) -- (0.419,-1.271) -- (0.375,-1.278) -- (0.331,-1.285) -- (0.283,-1.291) -- (0.234,-1.296) -- cycle;

  \end{scope}
  \coordinate (A) at (-1,0); \coordinate (B) at (1,0); \coordinate (C) at (0,2);
  \coordinate (F) at (0,0.5774); \coordinate (Cp) at (0,-1.7321);
  \fill[otNavy!8] (A) -- (B) -- (Cp) -- cycle;
  \draw[otNavy!60,line width=0.6pt] (A) -- (Cp) -- (B);
  \draw[otMuted,line width=0.5pt,dash pattern=on 2pt off 1.5pt] (F) -- (Cp);
  \draw[otInk,line width=1.1pt] (A) -- (B) -- (C) -- cycle;
  \draw[otGreenDark,line width=1.2pt] (F) -- (A) (F) -- (B) (F) -- (C);
  \foreach \p/\n/\a in {A/A/north east,B/B/north west,C/C/south}{%
    \fill[otInk] (\p) circle (1.6pt); \node[font=\small,anchor=\a] at (\p) {$\n$};}
  \filldraw[fill=otGreen,draw=otGreenDark] (F) circle (2pt);
  \node[font=\small,text=otGreenDark,anchor=west] at ([xshift=3pt,yshift=4pt]F) {$F$};
  \fill[otNavy] (Cp) circle (1.6pt);
  \node[font=\small,text=otNavy,anchor=north] at (Cp) {$C'$};
\end{tikzpicture}\\[0.5mm]
\legende{Lignes de niveau de $f$ (en gris bleuté) autour de son minimum $F$ ;
les segments $FA$, $FB$, $FC$ font entre eux des angles de $120^{\circ}$, et
$F$ est sur le segment $[CC']$, où $C'$ est le sommet du triangle
équilatéral extérieur construit sur $[AB]$.}
\end{center}

\partieExo{C}{Un exemple, et le retour de Napoléon}
On prend $A=(-1,0)$, $B=(1,0)$ et $C=(0,2)$.
\begin{questions}[resume]
  \item Vérifier que les angles du triangle sont inférieurs à
    $\frac{2\pi}{3}$. Par symétrie, chercher $F$ sur l'axe des ordonnées et
    montrer que $F=\bigl(0,\frac{1}{\sqrt3}\bigr)$. Calculer $f(F)$.
  \item Soit $C'$ le sommet du triangle équilatéral construit à l'extérieur
    sur $[AB]$. Montrer que $F\in[CC']$ et que $f(F)=CC'$ : la longueur
    minimale est la longueur commune $AA'=BB'=CC'$ du défi du chapitre 3.
  \item Interpréter la figure à l'aide du chapitre 20 : position du gradient
    par rapport aux lignes de niveau, et forme de celles-ci près de $F$ et
    loin du triangle.
\end{questions}
\end{defi}
